\documentclass[10pt]{article}
\usepackage[T1]{fontenc}
\usepackage[utf8]{inputenc}
\usepackage[french]{babel}
\usepackage{amsmath,amssymb,mathrsfs,mathpazo}
\usepackage{geometry}
\geometry{a4paper,margin=2cm}
\usepackage{graphicx,array,tabularx,tabularray,colortbl,multirow,multicol,booktabs,enumitem,pifont,ulem,fancybox,diagbox,cancel}
\usepackage{pstricks,pstricks-add,pst-plot,pst-tree,pst-node,pst-text,pst-func,pst-math,pst-bspline,pst-3dplot,pst-eucl,pst-solides3d}
\let\pscalculate\undefined
\usepackage{pst-bezier}
\usepackage{tikz,pgfplots,tkz-tab}
\usetikzlibrary{arrows,arrows.meta,patterns,calc,positioning,trees,decorations.markings,shapes.geometric,shapes.misc,shapes.arrows,intersections,decorations.pathmorphing,decorations.pathreplacing,backgrounds,shadows,babel}
\pgfplotsset{compat=1.18}
\usepgfplotslibrary{fillbetween}
\usepackage[np]{numprint}
\usepackage{listings,makecell,pas-tableur,twoopt}
\let\barre\undefined
\usepackage[pstricks]{tablvar}
\definecolor{darkgreen}{RGB}{0,100,0}
\providecommand{\red}{\color{red}}
\providecommand{\blue}{\color{blue}}
\providecommand{\green}{\color{green}}
\providecommand{\black}{\color{black}}
\providecommand{\white}{\color{white}}
\providecommand{\gray}{\color{gray}}
\providecommand{\cyan}{\color{cyan}}
\providecommand{\magenta}{\color{magenta}}
\providecommand{\yellow}{\color{yellow}}
\providecommand{\violet}{\color{violet}}
\providecommand{\orange}{\color{orange}}
\providecommand{\euro}{\texteuro}
\providecommand{\EUR}{\texteuro}
\providecommand{\ieme}{\textsuperscript{e}}
\providecommand{\degres}{\ensuremath{^\circ}}
\providecommand{\degree}{\ensuremath{^\circ}}
\providecommand{\vv}[1]{\overrightarrow{#1}}
\providecommand{\Coord}[1]{#1}
\providecommand{\decosix}{}
\providecommand{\starredbullet}{\ensuremath{\bullet}}
\pagestyle{empty}
\setlength{\parindent}{0pt}
\usepackage{adjustbox}\DeclareUnicodeCharacter{2236}{:}\DeclareUnicodeCharacter{2190}{\ensuremath{\leftarrow}}
\providecommand{\R}{\mathbb{R}}
\providecommand{\N}{\mathbb{N}}
\providecommand{\Z}{\mathbb{Z}}
\providecommand{\Q}{\mathbb{Q}}
\providecommand{\D}{\mathbb{D}}
\providecommand{\C}{\mathbb{C}}
\providecommand{\e}{{\mathrm{e}}}
\providecommand{\vect}[1]{\overrightarrow{#1}}
\providecommand{\vectt}[1]{\overrightarrow{\mathrm{#1}}}
\providecommand{\vvt}[1]{\overrightarrow{\mathrm{#1}}}
\providecommand{\np}[2][]{{#2~#1}}
\providecommand{\vv}[1]{{\overrightarrow{#1}}}
\renewcommand{\d}{{\mathrm{d}}}
\providecommand{\degres}{{^\circ}}
\providecommand{\degre}{{^\circ}}
\providecommand{\blue}{\color{blue}}
\providecommand{\red}{\color{red}}
\providecommand{\enskip}{\quad}
\providecommand{\llbracket}{\lbrack\!\lbrack}
\providecommand{\rrbracket}{\rbrack\!\rbrack}
\providecommand{\euro}{{\text{€}}}
\providecommand{\no}{{\text{n°}}}
\providecommand{\interval}[2]{{[#1;#2]}}
\providecommand{\intervFO}[2]{{[#1;#2[}}
\providecommand{\dotfill}{\ldots}

\providecommand{\ds}{\displaystyle}
\providecommand{\textding}[1]{#1}
\providecommand{\Oij}{\(\mathrm{O};\vec{\imath},\vec{\jmath}\)}
\providecommand{\Oijk}{\(\mathrm{O};\vec{\imath},\vec{\jmath},\vec{k}\)}
\providecommand{\Ouv}{\(\mathrm{O};\vec{u},\vec{v}\)}
\providecommand{\intervFF}[2]{\ensuremath{[#1;#2]}}
\renewcommand{\euro}{\ifmmode\text{\texteuro}\else\texteuro\fi}
\DeclareUnicodeCharacter{00B0}{\ifmmode^{\circ}\else\textdegree\fi}
\DeclareUnicodeCharacter{20AC}{\ifmmode\text{\texteuro}\else\texteuro\fi}
% Énoncés : Éducation nationale. Transcription et corrigés : APMEP, Denis Vergès et contributeurs. Découpage et classement : CoopMaths, Sébastien Lozano, Éric Elter et contributeurs. Adaptation de la présentation : Maths Explorer.
% CC BY-SA : https://coopmaths.fr/www/about/
% https://coopmaths.fr/alea/static/eam/2026/tex/eam_specifique_2026_06_metropole_2_cor.tex
\begin{document}
\noindent\textit{APMEP / CoopMaths. CC BY-SA. Présentation adaptée par Maths Explorer.}\par\medskip
\subsection*{A -- Premier placement.}
	\begin{enumerate}
		\item $a_{1} = a_0+200 = \np{20200}$ \quad et \quad $a_{2} = a_1 +200 = \np{20400}$.

		\item \begin{enumerate}
			\item Chaque année, le capital augmente de 200\euro{}, on a donc pour tout entier naturel $n$ :

			$a_{n+1}=a_n+200$.

			\item On reconnaît la relation de récurrence d'une suite arithmétique. La suite $(a_n)$ est arithmétique de raison $r=200$.
		\end{enumerate}

		\item La suite $(a_n)$ est arithmétique donc on a pour tout $n$ entier naturel :\quad $a_n=a_0+200n$.

		Soit ici,\quad $a_n=\np{20000} +200n$.

		\item On cherche $n$ tel que $a_n \geqslant  \np{22000}$.\hspace*{2cm}
		$\begin{aligned}[t] a_n \geqslant  \np{22000} & \iff \np{20000} +200n \geqslant  \np{22000}\\
		 & \iff 200n \geqslant  \np{2000}\\
		 & \iff n \geqslant \dfrac{\np{2000}}{200} \vspace*{2mm}\\
		 & \iff n \geqslant  10
		\end{aligned}$

		Ils devront attendre 10 ans soit à partir de 2035.
	\end{enumerate}

\subsection*{B -- Second placement.}
	\begin{enumerate}
		\item Le capital augmente de $2\%$, il est donc multiplié par 1,02.

			$b_1=b_0 \times 1,02 = \np{20000} \times 1,02=\np{20400}$.

		\item \begin{enumerate}
			\item En généralisant l'opération, chaque année, le capital est obtenu en multipliant le capital de l'année précédente par 1,02 :

			Pour tout $n$ entier naturel :\quad $b_{n+1}=1,02 b_n$.

			\item On reconnaît la relation de récurrence d'une suite géométrique.

			La suite $(b_n)$ est géométrique de raison $q=1,02$.
		\end{enumerate}

		\item La suite $(b_n)$ est géométrique donc on a, pour tout entier naturel $n$: \quad
		$b_n=b_0 \times q^n$

		soit $b_n= \np{20000} \times 1,02^n$

		\item Le tableau permet de dire que dès la 5\up{ème} année, ils auront le capital nécessaire à leur projet, soit en 2030.
	\end{enumerate}

\subsection*{C -- Bilan}
	Le deuxième placement est plus intéressant pour eux.

\end{document}
